\section{三类响应模型与检验路线}

通量限制、需求吸收和适应性反馈，可以产生不同的响应形态。本节列出三种待检验模型，并给出相应数据口径。选取哪种形式，取决于具体任务的机制与观察；介质分类为选择提供线索。

数字服务的一种情形，是在需求与分发条件相对稳定时，实际使用率向平台趋近。平台可以受到注意力、预算、供给和产品适配等因素影响。单位生成成本与实际使用率需要分别测量，快速扩张的候选并不自动转成使用。

这条曲线的形状可以写成，

\begin{equation}
Y_{\text{比特}}(t)=Y_{\infty}\left(1-e^{-t/\tau_{\text{趋近}}}\right)
\end{equation}

$Y_{\text{比特}}$表示选定服务的使用率或有效处理率，$Y_\infty$为同口径的平台值，$\tau_{\text{趋近}}$为时间常数。式中初始值取零，早期近似线性上升，随后以指数衰减的偏差接近平台。平台位置和趋近速度须由数据估计，需求或供给变化以后可采用新的参数或形式。

固定设备、空间与能源供应，可能形成单位时间的实物产出上限。投入接近这套设施的能力时，增加投入的收益可以递减；新增设备、改变工艺或扩建网络，则可能改变上限。检验可以分别比较设计、制造、施工与维护的投入和产出，辨认AI究竟改善了哪些环节。

这条曲线的两个特征可以写成一组条件，

\begin{equation}
0\le Y_{\text{原子}}(I)\le \Phi_{\max},\qquad \frac{\partial^{2}Y_{\text{原子}}}{\partial I^{2}}\le0\quad(\text{饱和近似})
\end{equation}

$Y_{\text{原子}}$为实物产出率，$\Phi_{\max}$为给定设施条件下的最大可持续产出率，两者采用相同单位；$I$为统一口径的投入。第一项表达能力上界，第二项另加了饱和区间的凹性近似。阶跃扩产、学习效应或互补投入可能不满足凹性，需按观察选择其他形式。

能量比较也须固定任务。擦除一个比特、断开一个化学键和抬升一个宏观物体，分别选择了不同操作单位；它们的绝对值不能决定各产业的增长率。对一项完整任务，应核算所需操作次数、实际耗散、设备能力与有效产出，再比较设计改进和执行改进带来的收益。

适应性反馈可以改变剂量与效果的关系。一个可检验情形是：超过某个强度以后，平均收益趋于平台，而结果的离散程度继续增加。也可能观察到收益持续改善、作用下降或迟发响应。实验需要同时记录强度、作用时间、群体差异与后续观察，依据数据辨认具体反馈。

把这条检验路线写成判据就是，

\begin{equation}
\left.\frac{d\langle Y_{\text{生命}}\rangle}{dP}\right|_{P>P_{c}}\approx0,\qquad \left.\frac{d\,\mathrm{Var}(Y_{\text{生命}})}{dP}\right|_{P>P_{c}}>0
\end{equation}

$Y_{\text{生命}}$为选定结果指标，$P$为采用固定单位的干预强度，$P_c$为待估计的转折位置。均值与方差采用各自单位。式中的两个条件共同描述这一组平台与风险扩大的观察，其他反馈形态需要另外的函数。比较时须保留群体、观察时长和重复测量信息。

这三个示例分别选取时间、投入和干预强度作自变量，因而不能直接比较图上的斜率。要比较同一项智能技术在不同环节的效果，应先统一任务、投入与结果口径，再分别控制过程时间和环境条件。图形的差别来自机制假设，是否对应实际领域差异，仍由这种比较检验。

\begin{figure}[ht]
\centering
% >>>>>>>>>> figs/fig-scaling-laws.tex >>>>>>>>>>
% 图 A-1  三条演化曲线的形状与各自的横轴
% 这张图要读的不是三条曲线的具体高度，而是它们各自配了不同的横轴：
% (a) 时间／累计投入，(b) 投入，(c) 干预强度。正文那一节的结论正是
% 「把三条曲线放进同一个坐标系比较进步速度没有意义，因为横轴不同」，
% 所以三个横轴标签是图里最不能省的三处文字，其余一律压到最短。
% 纵轴一律不给数值刻度：(a)(b) 是相对产出，(c) 是收益与方差，
% 标上刻度会诱导读者去做这张示意图支持不了的横向比较。
% (c) 里方差线穿过平均收益线是刻意的——「超过阈值之后平均收益不再增加
% 而不确定性继续上升」在这个穿越上才看得见。
% 子图编号 (a)(b)(c) 一律放图下方。曲线用 plot[smooth]，因为这三条
% 本来就是连续可导的函数，与 fig-proof-abundance 里用折线的情形不同。
\begin{tikzpicture}[x=1cm, y=1cm, font=\small,
  ax/.style={-{Stealth[length=2.0mm]}, ink, line width=0.6pt}
]
\useasboundingbox (-0.30,3.90) rectangle (13.35,-0.62);

% ==================== (a) 比特：单调趋近 ====================
\node[font=\scriptsize, text=slate, anchor=center] at (2.15,3.62) {单调趋近};

\draw[slate!55, line width=0.9pt, dashed] (0.50,2.90) -- (3.80,2.90);
\node[font=\scriptsize, text=slate, anchor=west] at (3.86,2.90) {$Y_{\infty}$};

\draw[deepblue, line width=1.9pt] plot[smooth] coordinates {
  (0.50,0.60) (0.78,1.11) (1.06,1.50) (1.34,1.81) (1.62,2.05)
  (1.90,2.24) (2.18,2.39) (2.46,2.51) (2.74,2.60) (3.02,2.67)
  (3.30,2.73) (3.58,2.77) (3.80,2.79)
};

\draw[ax] (0.50,0.60) -- (4.02,0.60);
\draw[ax] (0.50,0.60) -- (0.50,3.42);
\node[font=\scriptsize, text=ink, anchor=north] at (2.15,0.42) {时间};
\node[font=\scriptsize, text=ink, rotate=90, anchor=south] at (0.34,2.00) {产出};

% ==================== (b) 原子：受通量约束 ====================
\node[font=\scriptsize, text=slate, anchor=center] at (6.60,3.62) {受通量约束};

\fill[rust!10] (4.95,3.10) rectangle (8.25,3.40);
\draw[rust!70, line width=0.9pt, dashed] (4.95,3.10) -- (8.25,3.10);
\node[font=\scriptsize, text=rust!85, anchor=north east] at (8.25,3.06) {$\Phi_{\max}$};

\draw[terra, line width=1.9pt] plot[smooth] coordinates {
  (4.95,0.60) (5.28,1.56) (5.61,1.82) (5.94,2.01) (6.27,2.16)
  (6.60,2.29) (6.93,2.40) (7.26,2.50) (7.59,2.59) (7.92,2.67)
  (8.25,2.75)
};

\draw[ax] (4.95,0.60) -- (8.47,0.60);
\draw[ax] (4.95,0.60) -- (4.95,3.42);
\node[font=\scriptsize, text=ink, anchor=north] at (6.60,0.42) {投入};
\node[font=\scriptsize, text=ink, rotate=90, anchor=south] at (4.79,2.00) {产出};

% ==================== (c) 生命：反馈型 ====================
\node[font=\scriptsize, text=slate, anchor=center] at (11.05,3.62) {反馈型};

\draw[rust!70, line width=0.9pt, dashed] (11.215,0.60) -- (11.215,3.10);
\node[font=\scriptsize, text=rust!85, anchor=east] at (11.11,1.78) {$P_{c}$};

% 平均收益：升到阈值后转平
\draw[forest, line width=1.9pt] plot[smooth] coordinates {
  (9.40,0.60) (9.73,1.03) (10.06,1.42) (10.39,1.76) (10.72,2.06)
  (11.05,2.33) (11.22,2.45) (11.60,2.45) (12.15,2.45) (12.70,2.45)
};
% 不确定性：阈值前基本不动，阈值后迅速上升
\draw[rust, line width=1.9pt] plot[smooth] coordinates {
  (9.40,1.05) (9.95,1.06) (10.50,1.09) (11.05,1.14) (11.22,1.18)
  (11.60,1.55) (11.95,2.05) (12.30,2.58) (12.70,3.08)
};

\draw[forest, line width=1.5pt] (9.55,3.22) -- (9.98,3.22);
\node[font=\scriptsize, text=forest, anchor=west] at (10.05,3.22) {平均收益};
\draw[rust, line width=1.5pt] (9.55,2.92) -- (9.98,2.92);
\node[font=\scriptsize, text=rust, anchor=west] at (10.05,2.92) {不确定性};

\draw[ax] (9.40,0.60) -- (12.92,0.60);
\draw[ax] (9.40,0.60) -- (9.40,3.42);
\node[font=\scriptsize, text=ink, anchor=north] at (11.05,0.42) {干预强度};
\node[font=\scriptsize, text=ink, rotate=90, anchor=south] at (9.24,2.00) {各自归一化结果};

% ==================== 子图编号，一律放下面 ====================
\node[font=\small, text=ink] at (2.15,-0.28) {(a)};
\node[font=\small, text=ink] at (6.60,-0.28) {(b)};
\node[font=\small, text=ink] at (11.05,-0.28) {(c)};
\end{tikzpicture}

% <<<<<<<<<< figs/fig-scaling-laws.tex <<<<<<<<<<

\caption{三个响应模型的概念示意。(a) 使用率随时间趋近平台，(b) 给定设施内产出率随投入饱和，(c) 平均收益趋平而结果方差增加。各纵轴独立归一化，尤其(c)的均值与方差不共用物理单位；线条高度及交点不供跨指标比较。}
\label{fig:scaling-laws}
\end{figure}

把三条曲线并排列出，差别落在四列上。

\begin{table}[ht]
\centering
\caption{三条演化曲线的自变量、判据与可用的检验数据}
\label{tab:scaling-laws}
{\small
\begin{tabularx}{\textwidth}{@{}>{\RaggedRight\arraybackslash}p{2.45cm} >{\RaggedRight\arraybackslash}p{1.9cm} >{\RaggedRight\arraybackslash}p{4.0cm} >{\RaggedRight\arraybackslash}X@{}}
\toprule
曲线 & 自变量 & 判据 & 可用的检验数据 \\
\midrule
比特：单调趋近 & 时间 & $Y_{\text{比特}}\to Y_{\infty}$，平台与时间常数由任务数据估计 & 文本、图像、代码生成的单位成本时间序列，以及这些产出的实际使用率，而非产出量 \\
原子：受通量约束 & 投入 $I$ & $Y_{\text{原子}}\le\Phi_{\max}$，饱和区间另检验$\partial^{2}Y_{\text{原子}}/\partial I^{2}\le0$ & 建筑业劳动生产率、采矿与冶炼的能耗强度、半导体制造的单位投资；同一项目在设计阶段与施工阶段的工期压缩比 \\
生命：反馈型 & 干预强度 $P$ & 在 $P>P_{c}$ 上 $d\langle Y_{\text{生命}}\rangle/dP\approx0$，且 $d\,\mathrm{Var}(Y_{\text{生命}})/dP>0$ & 药物剂量、疗效与副作用的配对测量；编辑强度与脱靶率的配对测量 \\
\bottomrule
\end{tabularx}}
\end{table}

表 \ref{tab:scaling-laws} 的四列里，只有第三列可以脱离本书的术语使用：判据一栏写的是可以直接从数据里取到的条件，不要求读者先接受 $\eta$、$De$ 或三类介质的分类。这是这张表与全书其余部分的关系：分类是本书提出的，检验不是。三条曲线的判据彼此独立，任何一条被观测否定，对应的那条曲线就要改写，另外两条不受影响，因为它们各自的自变量与判据都是分开写出来的。

同一技术通常穿过多个环节，结果还受到相互作用影响。总产出未必是三条曲线的线性叠加。可行的起点，是选择边界稳定的任务，分别记录设计、执行、验证与维护的时间和成本，再用对照、分阶段部署或可比较项目检查机制。整体效果的解释，要在这些环节关系得到支持以后推进。
